% System Design
% Detailed design of storage engines, indexes, query pipeline,
% and benchmark isolation.

\subsection{Detailed Design}

\subsubsection{Storage Engines}

\paragraph{Knowledge Graph (Baseline)}

The KG engine stores facts as RDF-style triples using SPO/POS/OSP hash
indexes.  N-ary facts are decomposed via reification~\cite{manola2004rdf}:
an event node is created for each fact and the subject, predicate, and
each object become triples with the event node as subject.

Given a fact with arity $k$ and $m$ attributes, the KG generates
$k + m + 3$ triples (one for the event type, one each for subject and
predicate, one per object and attribute).  Reconstruction of the
original fact requires $k + m$ joins.

The KG implements the following indexes:

\begin{itemize}
    \item \textbf{SPO}: Subject $\rightarrow$ Predicate $\rightarrow$ Object.
    \item \textbf{POS}: Predicate $\rightarrow$ Object $\rightarrow$ Subject.
    \item \textbf{OSP}: Object $\rightarrow$ Subject $\rightarrow$ Predicate.
\end{itemize}

\paragraph{Sheaf Database (Primary)}

The SheafDB engine stores facts directly as sections over a poset of
contexts without decomposition.  Each fact is assigned to its maximal
context and restriction maps to sub-contexts are maintained.

The SheafDB implements the following indexes:

\begin{itemize}
    \item \textbf{Context Index}: Maps a context to the set of sections
        (facts) assigned to that context.
    \item \textbf{Stalk Index}: Maps an entity identifier to the set of
        contexts in which that entity appears.
    \item \textbf{Restriction Index}: Maps a pair of contexts
        $(c_{\text{broad}}, c_{\text{narrow}})$ to the corresponding
        restriction map.
    \item \textbf{Neighbourhood Index}: Maps a context to its immediate
        parents and children in the context poset.
    \item \textbf{Global Section Cache}: Caches the result of gluing all
        local sections into the global section (root context).
\end{itemize}

Context-scoped queries examine only the relevant context's sections,
avoiding the join overhead of the KG.  Global queries require gluing
compatible local sections, which scales with the number of contexts
rather than the number of triples.

\paragraph{Simplicial Database (Secondary)}

The SimplicialDB engine stores facts as simplices in a simplicial set.
An $n$-ary fact becomes an $n$-simplex whose vertices are the entities
and values involved.  Face maps correspond to restriction (dropping a
participant) and degeneracy maps to adding redundant information.

The SimplicialDB implements the following indexes:

\begin{itemize}
    \item \textbf{Simplex Index}: Maps a simplex identifier to its
        ordered list of vertices.
    \item \textbf{Incidence Matrix}: Records which $(n-1)$-simplices
        (faces) are contained in which $n$-simplices.
    \item \textbf{Boundary Operator}: Pre-computed matrix representing
        $\partial_n$ for each dimension $n$, satisfying
        $\partial_{n-1} \circ \partial_n = 0$.
    \item \textbf{Dimension Index}: Groups simplices by dimension for
        efficient traversal.
\end{itemize}

The simplicial representation provides a natural notion of
\textbf{homological consistency}: inconsistencies in the knowledge base
manifest as non-trivial homology groups.

\subsubsection{Query Pipeline}

All queries pass through the same pipeline regardless of engine:

\begin{enumerate}
    \item \textbf{Parsing}: A query string is parsed into a logical plan
        (a \texttt{Query} object with type, filters, and parameters).
    \item \textbf{Optimization}: The logical plan is transformed into a
        physical plan by the engine-specific optimizer, which applies
        cost-based rules (selectivity estimation, index selection,
        join ordering).
    \item \textbf{Execution}: The physical plan is executed against the
        storage engine.  The execution engine handles index lookups,
        data retrieval, and result construction.
    \item \textbf{Canonicalization}: Engine-native results are mapped
        back to \texttt{SemanticFact} instances via the canonical
        adapter.
    \item \textbf{Validation}: Results are checked against invariants.
    \item \textbf{Benchmark Collection}: Metrics (latency, memory, rows
        scanned) are recorded.
\end{enumerate}

\subsubsection{Benchmark Isolation}

The benchmark layer is strictly isolated from the storage engines:

\begin{itemize}
    \item It communicates only through the \texttt{DatabaseEngine}
        abstract interface.
    \item It never inspects engine-internal state.
    \item It compares results only through the canonical model.
    \item It has no knowledge of which engine is running.
\end{itemize}

This design prevents biased comparisons: the same benchmark code runs
identically against all engines.  Any difference in results is due to
the engine, not the benchmark framework.

\subsubsection{Configuration}

All engines share a common configuration system with three sources
(in increasing priority order):

\begin{enumerate}
    \item YAML configuration files (e.g., \texttt{configs/benchmark.yaml}).
    \item Environment variables (prefix \texttt{SFDB\_}).
    \item Programmatic overrides via the \texttt{Config} dataclass.
\end{enumerate}

The configuration system exposes settings for storage backend, index
parameters, logging, and benchmark parameters.  Every engine is
configurable through the same \texttt{Config} type.
